\section*{Chapter \ref{ch:rs:point-in-time-data-and-the-biases} --- Point-in-Time Data and the Biases}

\begin{solution}{exo:rs:point-in-time-data-and-the-biases:1}
The close of day $t$ is known only when the day ends, so trading at that same close uses it before it exists (unless the signal
is computed from prices a few minutes earlier and sent into the closing auction). The earliest honest trade is the next
session's open, or the next close in a close-to-close test.
\end{solution}

\begin{solution}{exo:rs:point-in-time-data-and-the-biases:2}
$w(r_S - r_F) = 0.15 \times (0.10 + 0.20) = 0.045$: 4.5 points a year.
\end{solution}

\begin{solution}{exo:rs:point-in-time-data-and-the-biases:3}
In January 2009 (four months after), $-403\,000$; in June 2009, after the February benchmark, $-321\,000$; in September 2011,
$-434\,000$.
\end{solution}

\begin{solution}{exo:rs:point-in-time-data-and-the-biases:4}
On stamps: the quotes stamped 09:59:59.5, 10:00:01.0 and 10:00:01.8. On arrival times (09:59:59.8, 10:00:01.3, 10:00:02.1): the
first quote for the first two decisions and the second for the third. The join on stamps lets the second decision use a quote
0.3 seconds before it existed at the firm.
\end{solution}

\begin{solution}{exo:rs:point-in-time-data-and-the-biases:5}
The first-release surprise is the final-data surprise minus the revision, so its variance is $90^2 + 106^2 = 19\,336$ (thousand
jobs squared), of which $11\,236$, or 58\%, is revision noise. A model fitted on final data sees a much cleaner signal than the one
the market traded.
\end{solution}

\begin{solution}{exo:rs:point-in-time-data-and-the-biases:6}
The history before a fund joins was never available in real time and is selected: only funds with good past returns join, so
their backfilled years raise the database average (\omterm{def:rs:point-in-time-data-and-the-biases:vintage}{backfill bias}). Remove it by using each fund's returns only from its date of
entry into the database (its \omterm{def:rs:point-in-time-data-and-the-biases:bitemporal}{knowledge time}), or by dropping a fixed initial period of every fund.
\end{solution}

\begin{solution}{exo:rs:point-in-time-data-and-the-biases:7}
Barrier $-1.2$: 3.5\% of names delisted a year, bias 7.0 points; barrier $-2.0$: 1.4\% and 3.8 points. Delisting return 0: 4.5
points; $-100\%$: 6.7 points. The barrier moves the bias more, because it sets how many names fail and how long they spend
declining; the delisting return affects only one month of each failure.
\end{solution}

\begin{solution}{exo:rs:point-in-time-data-and-the-biases:8}
Three biases at once. Survivorship: today's members exclude the companies that failed or shrank out of the index since 2000.
Index-membership look-ahead: many of today's members were added because they had risen, a fact unknown in 2000. \omterm{def:rs:point-in-time-data-and-the-biases:vintage}{Restatements}
and vendor corrections: today's fundamentals are not those published at the time. The test needs the index membership as of each
date, delisted companies with their delisting returns, and point-in-time fundamentals.
\end{solution}

\begin{solution}{pb:rs:point-in-time-data-and-the-biases:1}
\textbf{1.} $1 - e^{-1.6} = 79.8\%$. \textbf{2.} 2.2\% of names a year. \textbf{3.} 7.7\% a year. \textbf{4.} 12.9\%.
\textbf{5.} 5.2 points. \textbf{6.} 20.9\%. \textbf{7.} $r_F = -11.1\%$, $r_S = 12.6\%$ a year. \textbf{8.} $0.209 \times (12.6 +
11.1) = 5.0$ points; Question~5 averages month by month over a changing number of survivors (few early, many late) instead of
pooling name-months. \textbf{9.} 4.5 points. \textbf{10.} 6.7 points. \textbf{11.} 3.5\% a year, 7.0 points. \textbf{12.} 1.4\%
a year, 3.8 points. \textbf{13.} More names reach the barrier, so more name-months belong to failures ($w$ rises), while each
failure's decline is still steep. \textbf{14.} A strong profit: the fallen stocks that went on to fail are missing, and those that
remain are by construction the ones that recovered; the true result may have the opposite sign. \textbf{15.} From a database that
keeps delisted securities and their delisting returns, with universes as of each date. \textbf{16.} An estimate by delisting
reason, such as the $-30\%$ Shumway recovered for performance delistings; zero is optimistic, since performance delistings are
surprises that the last traded price does not reflect. \textbf{17.} Failing names are small when they fail, so a
capitalisation-weighted universe gives them little weight and suffers less; equal weighting, common in research, suffers most.
\textbf{18.} \emph{Named result:} the survivors-only universe overstates the equal-weighted return by 5.2 points a year when 2.2\%
of names delist each year, and the delisting month explains 0.7 of them (the bias is still 4.5 points with a delisting return of
zero). \textbf{19.} \omterm{def:rs:point-in-time-data-and-the-biases:vintage}{Backfill bias} in fund and vendor databases, and testing on today's index members. \textbf{20.} It lists the
companies that will survive, which no one knew at the start of the sample.
\end{solution}

\begin{solution}{iq:rs:point-in-time-data-and-the-biases:1}
Studying only securities that exist at the end of the sample. In a stock-selection backtest built on today's listed companies,
the failures are missing, so strategies that buy cheap, distressed or small stocks look far better than they were, and average
returns are inflated by several points a year in equal-weighted universes.
\iqlookfor{that the missing names are the losers, and the fix: a universe as of each date with delistings.}
\end{solution}

\begin{solution}{iq:rs:point-in-time-data-and-the-biases:2}
The estimate available at each decision date: the vintage published by then (first or later release), not today's revised
figure; the model must also be fitted on those real-time values, since it will be fed them live.
\iqlookfor{real-time vintages; the distinction between valid and \omterm{def:rs:point-in-time-data-and-the-biases:bitemporal}{knowledge time}.}
\end{solution}

\begin{solution}{iq:rs:point-in-time-data-and-the-biases:3}
Key (company, field, fiscal period, \omterm{def:rs:point-in-time-data-and-the-biases:bitemporal}{knowledge time}) with the value, where \omterm{def:rs:point-in-time-data-and-the-biases:bitemporal}{knowledge time} is the filing or vendor delivery time;
never overwrite. A query takes (company, field, period, as-of time) and returns the latest row known at that time; a snapshot
query returns every period as known then.
\iqlookfor{two time columns, append-only storage, as-of semantics.}
\end{solution}

\begin{solution}{iq:rs:point-in-time-data-and-the-biases:4}
Perturbation: for each decision time, replace every input value with a \omterm{def:rs:point-in-time-data-and-the-biases:bitemporal}{knowledge time} after it by noise and check that the
features at that decision are unchanged. Also a guard in the data access layer that raises on any later \omterm{def:rs:point-in-time-data-and-the-biases:bitemporal}{knowledge time}, and
a comparison of the backtest run on stored vintages with the same run on today's data (a large difference flags a sensitivity to
revisions).
\iqlookfor{an automated test, not a code review.}
\end{solution}

\begin{solution}{iq:rs:point-in-time-data-and-the-biases:5}
Survivorship (dead funds dropped), backfill (funds join with their good past), self-selection (funds report when they do well
and stop when they do badly), and stale or smoothed marks for illiquid holdings, which lower measured volatility.
\iqlookfor{several distinct biases and their direction.}
\end{solution}

\begin{solution}{iq:rs:point-in-time-data-and-the-biases:6}
Pooled over name-months, the universe average is $(1 - w)r_S + wr_F$; the survivors' average is $r_S$; the difference is $w(r_S
- r_F)$. In practice the gap $r_S - r_F$ is dominated by the months failing names spend declining, not the delisting month: in the
chapter's simulation, 4.5 of the 5.2 points remain with a delisting return of zero.
\iqlookfor{the identity, and where the gap comes from.}
\end{solution}
